% Generated by report/tools/tables.py from the result files of the repository. Do not edit.
\begin{table}[tbp]
\centering
\footnotesize
\setlength{\tabcolsep}{4pt}
\caption[Accord des doubles lectures, avec intervalles]{Accord des doubles lectures à l'aveugle, toutes paires réunies, avec ses intervalles. 30 articles relus : G et G2, les mêmes 10 articles de 2026 avec les deux grilles ; R2, 10 articles ajoutés (5 de 2018, 5 de 2026) ; R3, 10 articles de la trajectoire. Accord : décisions identiques sur les articles où les deux lecteurs ont répondu ; intervalle de Wilson à 95~\% ; kappa de Cohen avec un intervalle bootstrap percentile (vide quand un lecteur ne varie pas). Sur les quatre critères qui portent un verdict (Q5, Y9, X7, X2), 120 décisions identiques sur 120. \textbf{En gras}, les critères dont le kappa reste sous 0,70. Avec dix articles par paire, 9 sur 10 a pour intervalle [60\,;\,98]~\% : l'accord d'une paire seule borne mal la fiabilité. Source : \file{study/survey\_extra.json} (\file{study/survey\_extra.py}).}
\label{tab:accord-intervalles}
\begin{tabularx}{\linewidth}{@{}l>{\raggedright\arraybackslash}Xrrr@{}}
\toprule
Code & Critère & accord & \% [Wilson 95~\%] & $\kappa$ [IC 95~\%] \\
\midrule
Q2 & Réseaux de neurones employés & 30/30 & 100 [89\,;\,100] & 1,00 [1,00\,;\,1,00] \\
Q3 & Tous les hyperparamètres de la méthode & 28/30 & 93 [79\,;\,98] & 0,86 [0,65\,;\,1,00] \\
Q4 & Tous les hyperparamètres des références & 30/30 & 100 [89\,;\,100] & 1,00 [1,00\,;\,1,00] \\
Q5 & Code en lien & 30/30 & 100 [89\,;\,100] & 1,00 [1,00\,;\,1,00] \\
Q6 & Choix des hyperparamètres décrit & 27/30 & 90 [74\,;\,97] & 0,75 [0,44\,;\,1,00] \\
Q7 & Évaluation sur conditions tenues à part & 27/30 & 90 [74\,;\,97] & 0,79 [0,55\,;\,1,00] \\
Q8 & Test statistique appliqué & 29/30 & 97 [83\,;\,99] & 0,84 [0,00\,;\,1,00] \\
\midrule
X2 & Incertitude rapportée & 30/30 & 100 [89\,;\,100] & 1,00 [1,00\,;\,1,00] \\
X3 & Seuls les meilleurs essais gardés & 29/30 & 97 [83\,;\,99] & 0,78 [0,00\,;\,1,00] \\
X4 & Références réglées avec le même soin & 29/30 & 97 [83\,;\,99] & 0,89 [0,62\,;\,1,00] \\
X7 & Calcul chiffré & 30/30 & 100 [89\,;\,100] & 1,00 [1,00\,;\,1,00] \\
\midrule
\textbf{Y1} & \textbf{Mesure de performance définie} & \textbf{26/30} & \textbf{87 [70\,;\,95]} & \textbf{0,69 [0,38\,;\,0,93]} \\
Y2 & Barres d'erreur définies & 24/24 & 100 [86\,;\,100] & 1,00 [1,00\,;\,1,00] \\
Y3 & Plusieurs tâches (au moins 2) & 30/30 & 100 [89\,;\,100] & 1,00 [1,00\,;\,1,00] \\
Y4 & Une tâche perdue montrée & 25/27 & 93 [77\,;\,98] & 0,81 [0,47\,;\,1,00] \\
\textbf{Y5} & \textbf{Origine du code des références} & \textbf{25/30} & \textbf{83 [66\,;\,93]} & \textbf{0,65 [0,33\,;\,0,92]} \\
\textbf{Y6} & \textbf{Exclusions dites} & \textbf{25/30} & \textbf{83 [66\,;\,93]} & \textbf{0,62 [0,27\,;\,0,89]} \\
Y7 & Version des bancs d'essai & 30/30 & 100 [89\,;\,100] & 1,00 [1,00\,;\,1,00] \\
Y8 & Test de généralisation & 28/30 & 93 [79\,;\,98] & 0,84 [0,59\,;\,1,00] \\
Y9 & Sensibilité à un hyperparamètre & 30/30 & 100 [89\,;\,100] & 1,00 [1,00\,;\,1,00] \\
Y10 & Versions des bibliothèques & 30/30 & 100 [89\,;\,100] & 1,00 [1,00\,;\,1,00] \\
Y11 & Type d'incertitude nommé & 24/24 & 100 [86\,;\,100] & 1,00 [1,00\,;\,1,00] \\
Y12 & Outil statistique postérieur à 2018 & 30/30 & 100 [89\,;\,100] & 1,00 [1,00\,;\,1,00] \\
\bottomrule
\end{tabularx}
\end{table}
